\IfFileExists{stacks-project.cls}{%
\documentclass{stacks-project}
}{%
\documentclass{amsart}
}

% For dealing with references we use the comment environment
\usepackage{verbatim}
\newenvironment{reference}{\comment}{\endcomment}
%\newenvironment{reference}{}{}
\newenvironment{slogan}{\comment}{\endcomment}
\newenvironment{history}{\comment}{\endcomment}

% For commutative diagrams we use Xy-pic
\usepackage[all]{xy}

% We use 2cell for 2-commutative diagrams.
\xyoption{2cell}
\UseAllTwocells

% We use multicol for the list of chapters between chapters
\usepackage{multicol}

% This is generally recommended for better output
\usepackage{lmodern}
\usepackage[T1]{fontenc}

% For cross-file-references
\usepackage{xr-hyper}

% Package for hypertext links:
\usepackage{hyperref}

% For any local file, say "hello.tex" you want to link to please
% use \externaldocument[hello-]{hello}
\externaldocument[introduction-]{introduction}
\externaldocument[conventions-]{conventions}
\externaldocument[sets-]{sets}
\externaldocument[categories-]{categories}
\externaldocument[topology-]{topology}
\externaldocument[sheaves-]{sheaves}
\externaldocument[sites-]{sites}
\externaldocument[stacks-]{stacks}
\externaldocument[fields-]{fields}
\externaldocument[algebra-]{algebra}
\externaldocument[brauer-]{brauer}
\externaldocument[homology-]{homology}
\externaldocument[derived-]{derived}
\externaldocument[simplicial-]{simplicial}
\externaldocument[more-algebra-]{more-algebra}
\externaldocument[smoothing-]{smoothing}
\externaldocument[modules-]{modules}
\externaldocument[sites-modules-]{sites-modules}
\externaldocument[injectives-]{injectives}
\externaldocument[cohomology-]{cohomology}
\externaldocument[sites-cohomology-]{sites-cohomology}
\externaldocument[dga-]{dga}
\externaldocument[dpa-]{dpa}
\externaldocument[sdga-]{sdga}
\externaldocument[hypercovering-]{hypercovering}
\externaldocument[schemes-]{schemes}
\externaldocument[constructions-]{constructions}
\externaldocument[properties-]{properties}
\externaldocument[morphisms-]{morphisms}
\externaldocument[coherent-]{coherent}
\externaldocument[divisors-]{divisors}
\externaldocument[limits-]{limits}
\externaldocument[varieties-]{varieties}
\externaldocument[topologies-]{topologies}
\externaldocument[descent-]{descent}
\externaldocument[perfect-]{perfect}
\externaldocument[more-morphisms-]{more-morphisms}
\externaldocument[flat-]{flat}
\externaldocument[groupoids-]{groupoids}
\externaldocument[more-groupoids-]{more-groupoids}
\externaldocument[etale-]{etale}
\externaldocument[chow-]{chow}
\externaldocument[intersection-]{intersection}
\externaldocument[pic-]{pic}
\externaldocument[weil-]{weil}
\externaldocument[adequate-]{adequate}
\externaldocument[dualizing-]{dualizing}
\externaldocument[duality-]{duality}
\externaldocument[discriminant-]{discriminant}
\externaldocument[derham-]{derham}
\externaldocument[local-cohomology-]{local-cohomology}
\externaldocument[algebraization-]{algebraization}
\externaldocument[curves-]{curves}
\externaldocument[resolve-]{resolve}
\externaldocument[models-]{models}
\externaldocument[functors-]{functors}
\externaldocument[equiv-]{equiv}
\externaldocument[pione-]{pione}
\externaldocument[etale-cohomology-]{etale-cohomology}
\externaldocument[proetale-]{proetale}
\externaldocument[relative-cycles-]{relative-cycles}
\externaldocument[more-etale-]{more-etale}
\externaldocument[trace-]{trace}
\externaldocument[crystalline-]{crystalline}
\externaldocument[spaces-]{spaces}
\externaldocument[spaces-properties-]{spaces-properties}
\externaldocument[spaces-morphisms-]{spaces-morphisms}
\externaldocument[decent-spaces-]{decent-spaces}
\externaldocument[spaces-cohomology-]{spaces-cohomology}
\externaldocument[spaces-limits-]{spaces-limits}
\externaldocument[spaces-divisors-]{spaces-divisors}
\externaldocument[spaces-over-fields-]{spaces-over-fields}
\externaldocument[spaces-topologies-]{spaces-topologies}
\externaldocument[spaces-descent-]{spaces-descent}
\externaldocument[spaces-perfect-]{spaces-perfect}
\externaldocument[spaces-more-morphisms-]{spaces-more-morphisms}
\externaldocument[spaces-flat-]{spaces-flat}
\externaldocument[spaces-groupoids-]{spaces-groupoids}
\externaldocument[spaces-more-groupoids-]{spaces-more-groupoids}
\externaldocument[bootstrap-]{bootstrap}
\externaldocument[spaces-pushouts-]{spaces-pushouts}
\externaldocument[spaces-chow-]{spaces-chow}
\externaldocument[groupoids-quotients-]{groupoids-quotients}
\externaldocument[spaces-more-cohomology-]{spaces-more-cohomology}
\externaldocument[spaces-simplicial-]{spaces-simplicial}
\externaldocument[spaces-duality-]{spaces-duality}
\externaldocument[formal-spaces-]{formal-spaces}
\externaldocument[restricted-]{restricted}
\externaldocument[spaces-resolve-]{spaces-resolve}
\externaldocument[formal-defos-]{formal-defos}
\externaldocument[defos-]{defos}
\externaldocument[cotangent-]{cotangent}
\externaldocument[examples-defos-]{examples-defos}
\externaldocument[algebraic-]{algebraic}
\externaldocument[examples-stacks-]{examples-stacks}
\externaldocument[stacks-sheaves-]{stacks-sheaves}
\externaldocument[criteria-]{criteria}
\externaldocument[artin-]{artin}
\externaldocument[quot-]{quot}
\externaldocument[stacks-properties-]{stacks-properties}
\externaldocument[stacks-morphisms-]{stacks-morphisms}
\externaldocument[stacks-limits-]{stacks-limits}
\externaldocument[stacks-cohomology-]{stacks-cohomology}
\externaldocument[stacks-perfect-]{stacks-perfect}
\externaldocument[stacks-introduction-]{stacks-introduction}
\externaldocument[stacks-more-morphisms-]{stacks-more-morphisms}
\externaldocument[stacks-geometry-]{stacks-geometry}
\externaldocument[moduli-]{moduli}
\externaldocument[moduli-curves-]{moduli-curves}
\externaldocument[examples-]{examples}
\externaldocument[exercises-]{exercises}
\externaldocument[guide-]{guide}
\externaldocument[desirables-]{desirables}
\externaldocument[coding-]{coding}
\externaldocument[obsolete-]{obsolete}
\externaldocument[fdl-]{fdl}
\externaldocument[index-]{index}

% Theorem environments.
%
\theoremstyle{plain}
\newtheorem{theorem}[subsection]{Theorem}
\newtheorem{proposition}[subsection]{Proposition}
\newtheorem{lemma}[subsection]{Lemma}

\theoremstyle{definition}
\newtheorem{definition}[subsection]{Definition}
\newtheorem{example}[subsection]{Example}
\newtheorem{exercise}[subsection]{Exercise}
\newtheorem{situation}[subsection]{Situation}

\theoremstyle{remark}
\newtheorem{remark}[subsection]{Remark}
\newtheorem{remarks}[subsection]{Remarks}

\numberwithin{equation}{subsection}

% Macros
%
\def\lim{\mathop{\mathrm{lim}}\nolimits}
\def\colim{\mathop{\mathrm{colim}}\nolimits}
\def\Spec{\mathop{\mathrm{Spec}}}
\def\Hom{\mathop{\mathrm{Hom}}\nolimits}
\def\Ext{\mathop{\mathrm{Ext}}\nolimits}
\def\SheafHom{\mathop{\mathcal{H}\!\mathit{om}}\nolimits}
\def\SheafExt{\mathop{\mathcal{E}\!\mathit{xt}}\nolimits}
\def\Sch{\mathit{Sch}}
\def\Mor{\mathop{\mathrm{Mor}}\nolimits}
\def\Ob{\mathop{\mathrm{Ob}}\nolimits}
\def\Sh{\mathop{\mathit{Sh}}\nolimits}
\def\NL{\mathop{N\!L}\nolimits}
\def\CH{\mathop{\mathrm{CH}}\nolimits}
\def\proetale{{pro\text{-}\acute{e}tale}}
\def\etale{{\acute{e}tale}}
\def\QCoh{\mathit{QCoh}}
\def\Ker{\mathop{\mathrm{Ker}}}
\def\Im{\mathop{\mathrm{Im}}}
\def\Coker{\mathop{\mathrm{Coker}}}
\def\Coim{\mathop{\mathrm{Coim}}}

% Boxtimes
%
\DeclareMathSymbol{\boxtimes}{\mathbin}{AMSa}{"02}

%
% Macros for moduli stacks/spaces
%
\def\QCohstack{\mathcal{QC}\!\mathit{oh}}
\def\Cohstack{\mathcal{C}\!\mathit{oh}}
\def\Spacesstack{\mathcal{S}\!\mathit{paces}}
\def\Quotfunctor{\mathrm{Quot}}
\def\Hilbfunctor{\mathrm{Hilb}}
\def\Curvesstack{\mathcal{C}\!\mathit{urves}}
\def\Polarizedstack{\mathcal{P}\!\mathit{olarized}}
\def\Complexesstack{\mathcal{C}\!\mathit{omplexes}}
% \Pic is the operator that assigns to X its picard group, usage \Pic(X)
% \Picardstack_{X/B} denotes the Picard stack of X over B
% \Picardfunctor_{X/B} denotes the Picard functor of X over B
\def\Pic{\mathop{\mathrm{Pic}}\nolimits}
\def\Picardstack{\mathcal{P}\!\mathit{ic}}
\def\Picardfunctor{\mathrm{Pic}}
\def\Deformationcategory{\mathcal{D}\!\mathit{ef}}
\usepackage{cleveref}
\crefname{originallemma}{original lemma}{original lemmas}
\Crefname{originallemma}{Original lemma}{Original lemmas}
\makeatletter
\newlabel{lemma-2-product-stacks}{{026G}{source}{}{sourcepointer-1}{}}
\newlabel{lemma-2-product-stacks@cref}{{[originallemma][0][]026G}{[1][0][]source}{}{}{}}
\makeatother
\crefname{originallemma}{original lemma}{original lemmas}
\Crefname{originallemma}{Original lemma}{Original lemmas}
\makeatletter
\newlabel{lemma-characterize-essentially-surjective-when-ff}{{046N}{source}{}{sourcepointer-2}{}}
\newlabel{lemma-characterize-essentially-surjective-when-ff@cref}{{[originallemma][0][]046N}{[1][0][]source}{}{}{}}
\makeatother
\crefname{originallemma}{original lemma}{original lemmas}
\Crefname{originallemma}{Original lemma}{Original lemmas}
\makeatletter
\newlabel{lemma-characterize-ff}{{04WQ}{source}{}{sourcepointer-3}{}}
\newlabel{lemma-characterize-ff@cref}{{[originallemma][0][]04WQ}{[1][0][]source}{}{}{}}
\makeatother
\crefname{originallemma}{original lemma}{original lemmas}
\Crefname{originallemma}{Original lemma}{Original lemmas}
\makeatletter
\newlabel{lemma-stackify}{{02ZN}{source}{}{sourcepointer-4}{}}
\newlabel{lemma-stackify@cref}{{[originallemma][0][]02ZN}{[1][0][]source}{}{}{}}
\makeatother
\crefname{originallemma}{original lemma}{original lemmas}
\Crefname{originallemma}{Original lemma}{Original lemmas}
\makeatletter
\newlabel{lemma-stackify-universal-property}{{0435}{source}{}{sourcepointer-5}{}}
\newlabel{lemma-stackify-universal-property@cref}{{[originallemma][0][]0435}{[1][0][]source}{}{}{}}
\makeatother
\crefname{originalreference}{original reference}{original references}
\Crefname{originalreference}{Original reference}{Original references}
\makeatletter
\newlabel{sites-section-glueing-sheaves}{{04TP}{source}{}{sourcepointer-6}{}}
\newlabel{sites-section-glueing-sheaves@cref}{{[originalreference][0][]04TP}{[1][0][]source}{}{}{}}
\makeatother
\title{Stackification commutes with 2-fibre products of fibred categories}
\author{The Stacks Project Authors}
\date{stacks-project:stacks.tex:local-snapshot-20261012; GNU FDL 1.2 or later}
\begin{document}
\maketitle
\noindent\textit{Editorial note.} No omission markers in-span; proof complete as written.\par\medskip
\noindent\textit{Allowed standard prerequisites (editorial).}\begin{itemize}
\item Universal property of stackification. lemma-stackify-universal-property (same chapter, named). Extends the canonical map from the fibre product to the fibre product of the stackifications, giving h : W -> X' x\_\{Y'\} Z'.
\item Stack criterion for a morphism. lemma-characterize-ff (same chapter, named). Reduces equivalence of h to isomorphisms of Mor sheaves plus local essential surjectivity.
\item Essential surjectivity criterion. lemma-characterize-essentially-surjective-when-ff (same chapter, named). Supplies the local-essential-image condition paired with the Mor-sheaf isomorphisms.
\item Construction of the stackification. lemma-stackify (same chapter, named). Identifies Mor sheaves of W locally with sheafified fibre products of Mor presheaves.
\end{itemize}\medskip
\ifdefined\hypertarget\hypertarget{sourcepointer-1}{}\fi
\noindent\textit{Source reference.} Original lemma 026G: sources/stacks-stacks/pinned-original-source.tex, source lines 939--939. Navigation pointer only: the counted proof cites this statement by \textbackslash{}ref as a standing input; the deduction the counted proof has to make is carried out inside the counted span, and no part of that deduction is delegated to this pointer.\par\medskip
\ifdefined\hypertarget\hypertarget{sourcepointer-2}{}\fi
\noindent\textit{Source reference.} Original lemma 046N: sources/stacks-stacks/pinned-original-source.tex, source lines 1062--1062. Navigation pointer only: the counted proof cites this statement by \textbackslash{}ref as a standing input; the deduction the counted proof has to make is carried out inside the counted span, and no part of that deduction is delegated to this pointer.\par\medskip
\ifdefined\hypertarget\hypertarget{sourcepointer-3}{}\fi
\noindent\textit{Source reference.} Original lemma 04WQ: sources/stacks-stacks/pinned-original-source.tex, source lines 1009--1009. Navigation pointer only: the counted proof cites this statement by \textbackslash{}ref as a standing input; the deduction the counted proof has to make is carried out inside the counted span, and no part of that deduction is delegated to this pointer.\par\medskip
\ifdefined\hypertarget\hypertarget{sourcepointer-4}{}\fi
\noindent\textit{Source reference.} Original lemma 02ZN: sources/stacks-stacks/pinned-original-source.tex, source lines 1697--1697. Navigation pointer only: the counted proof cites this statement by \textbackslash{}ref as a standing input; the deduction the counted proof has to make is carried out inside the counted span, and no part of that deduction is delegated to this pointer.\par\medskip
\ifdefined\hypertarget\hypertarget{sourcepointer-5}{}\fi
\noindent\textit{Source reference.} Original lemma 0435: sources/stacks-stacks/pinned-original-source.tex, source lines 1948--1948. Navigation pointer only: the counted proof cites this statement by \textbackslash{}ref as a standing input; the deduction the counted proof has to make is carried out inside the counted span, and no part of that deduction is delegated to this pointer.\par\medskip
\ifdefined\hypertarget\hypertarget{sourcepointer-6}{}\fi
\noindent\textit{Source reference.} Original reference 04TP: focused-seeds/geometry/S9-stk-stackification-fibre-product-fibred-categories/ref-tags.txt, source lines 5--5. Navigation pointer only: the counted proof cites this statement by \textbackslash{}ref as a standing input; the deduction the counted proof has to make is carried out inside the counted span, and no part of that deduction is delegated to this pointer.\par\medskip
\begin{lemma}
\label{lemma-stackification-fibre-product-fibred-categories}
Let $\mathcal{C}$ be a site.
Let $f : \mathcal{X} \to \mathcal{Y}$ and $g : \mathcal{Z} \to \mathcal{Y}$
be morphisms of fibred categories over $\mathcal{C}$.
In this case the stackification of the $2$-fibre product is the $2$-fibre
product of the stackifications.
\end{lemma}

\begin{proof}
Let us denote $\mathcal{X}', \mathcal{Y}', \mathcal{Z}'$ the stackifications
and $\mathcal{W}$ the stackification of
$\mathcal{X} \times_\mathcal{Y} \mathcal{Z}$. By construction of $2$-fibre
products there is a canonical $1$-morphism
$\mathcal{X} \times_\mathcal{Y} \mathcal{Z} \to
\mathcal{X}' \times_{\mathcal{Y}'} \mathcal{Z}'$.
As the second $2$-fibre product is a stack (see
Lemma \ref{lemma-2-product-stacks})
this $1$-morphism induces a $1$-morphism
$h : \mathcal{W} \to \mathcal{X}' \times_{\mathcal{Y}'} \mathcal{Z}'$
by the universal property of stackification, see
Lemma \ref{lemma-stackify-universal-property}.
Now $h$ is a morphism of stacks, and we may check that it is an
equivalence using
Lemmas \ref{lemma-characterize-ff} and
\ref{lemma-characterize-essentially-surjective-when-ff}.

\medskip\noindent
Thus we first prove that $h$ induces isomorphisms of $\mathit{Mor}$-sheaves.
Let $\xi, \xi'$ be objects of $\mathcal{W}$ over
$U \in \Ob(\mathcal{C})$. We want to show that
$$
h : \mathit{Mor}(\xi, \xi') \longrightarrow \mathit{Mor}(h(\xi), h(\xi'))
$$
is an isomorphism. To do this we may work locally on $U$ (see
Sites, Section \ref{sites-section-glueing-sheaves}).
Hence by construction of $\mathcal{W}$ (see
Lemma \ref{lemma-stackify})
we may assume that $\xi, \xi'$
actually come from objects $(x, z, \alpha)$ and $(x', z', \alpha')$
of $\mathcal{X} \times_\mathcal{Y} \mathcal{Z}$ over $U$.
By the same lemma once more we see that in this case
$\mathit{Mor}(\xi, \xi')$ is the sheafification of
$$
V/U \longmapsto
\Mor_{\mathcal{X}_V}(x|_V, x'|_V)
\times_{\Mor_{\mathcal{Y}_V}(f(x)|_V, f(x')|_V)}
\Mor_{\mathcal{Z}_V}(z|_V, z'|_V)
$$
and that $\mathit{Mor}(h(\xi), h(\xi'))$ is equal to the fibre product
$$
\mathit{Mor}(i(x), i(x'))
\times_{\mathit{Mor}(j(f(x)), j(f(x'))}
\mathit{Mor}(k(z), k(z'))
$$
where $i : \mathcal{X} \to \mathcal{X}'$,
$j : \mathcal{Y} \to \mathcal{Y}'$, and
$k : \mathcal{Z} \to \mathcal{Z}'$ are the canonical functors.
Thus the first displayed map of this paragraph is an isomorphism as
sheafification is exact (and hence the sheafification of a fibre product
of presheaves is the fibre product of the sheafifications).

\medskip\noindent
Finally, we have to check that any object of
$\mathcal{X}' \times_{\mathcal{Y}'} \mathcal{Z}'$
over $U$ is locally on $U$ in the essential image of $h$.
Write such an object as a triple $(x', z', \alpha)$.
Then $x'$ locally comes from an object of $\mathcal{X}$,
$z'$ locally comes from an object of $\mathcal{Z}$, and
having made suitable replacements for $x'$, $z'$ the morphism
$\alpha$ of $\mathcal{Y}'_U$ locally comes from a morphism of
$\mathcal{Y}$. In other words, we have shown that any object of
$\mathcal{X}' \times_{\mathcal{Y}'} \mathcal{Z}'$
over $U$ is locally on $U$ in the essential image of
$\mathcal{X} \times_\mathcal{Y} \mathcal{Z} \to
\mathcal{X}' \times_{\mathcal{Y}'} \mathcal{Z}'$, hence
a fortiori it is locally in the essential image of $h$.
\end{proof}

\end{document}
